\section{Conclusion and Future Work}
\label{ch:conclusion}

\subsection{Conclusions}

This project built a walk-forward back-testing engine for monthly SIPs in Nifty~50, gold and a
liquid fund, applied it to Indian data from 2000 to 2026, and turned the final rule into a public
planning tool. The main conclusions are:

\begin{enumerate}
  \item \textbf{Diversification across these three assets works.} Their correlations were low or
        negative throughout, and in every stress period studied the mixes lost far less than equity
        alone.
  \item \textbf{The risk-free benchmark matters.} With a cash-like liquid fund available, optimisers
        that measure risk against 0\% simply hold cash. Measuring returns above the liquid fund and
        treating it as cash, following \citet{tobin1958} and \citet{sharpe1994}, raised the
        optimised SIP's out-of-sample return from 11.1\% to 20.3\% a year.
  \item \textbf{Estimation error dominates.} The strategy that directly maximised the investor's
        utility never achieved the highest utility. Rules that depend less on estimated average
        returns, such as equal weights, fixed mixes and the risk-parity blend, held up better, in
        line with \citet{michaud1989}, \citet{chopra1993} and \citet{demiguel2009}.
  \item \textbf{The best strategy depends on the investor and the period.} Measured by utility, the
        ranking changed with risk aversion in 2010--2019 and with the gold boom in 2020--2025. No
        rule should be presented as universally best.
  \item \textbf{A personal cushion makes the advice personal.} The Personal SIP gives a conservative
        investor a much smaller worst fall ($-3.4\%$ against $-12.2\%$), at the cost of lower
        returns in strong markets. Its cushions, chosen on 2010--2019, proved too cautious for
        2020--2025.
  \item \textbf{The rule can be used in practice.} The planner turns the back-tested rule into a
        two-minute monthly routine, with results verified against the research code.
\end{enumerate}

\subsection{Limitations and future work}

The data before 2020 uses a constant 1.3\% Nifty dividend yield, international gold without Indian
import duty, and a T-bill index rather than a real liquid fund. The test period is only six years.
The values $\lambda = 2, 4, 8$ are assumed rather than measured, variance penalises gains as much as
losses, and taxes and exit loads are ignored. Natural extensions are: setting $\lambda$ from a risk
questionnaire; adding a drawdown penalty to $U$; shrinking estimated returns to reduce estimation
error; fetching fund prices automatically each month; and including taxes when deciding whether to
rebalance.
